%=====================================================================
%  Chapter 9 — General Conclusion and Perspectives
%=====================================================================
\chapter{General Conclusion and Perspectives}
\label{chap:conclusion}

\section{Summary of the Work}
This project set out to design and build, from scratch, a child-sized
18-DOF humanoid robot, \textbf{Abbes}, and --- with a software-engineering
emphasis --- to develop the complete software that makes it stand, balance,
move and respond to commands. The work spanned the entire stack. The
\emph{mechanical} design derived a human-like, eighteen-degree-of-freedom
morphology from anthropometric data and realised it in 3D-printed parts.
The \emph{electronics} were organised around an ESP32 microcontroller, a
PCA9685 PWM driver and an MPU6050 inertial sensor. The \emph{firmware}
turned that microcontroller into a real-time, self-reconnecting micro-ROS
node, making the on-robot controller a first-class participant in the
ROS\,2 graph. On top of it, a modular \emph{ROS\,2 software architecture}
separated the robot model, the simulation bring-up, the algorithms and the
interfaces into independent packages, and used \texttt{ros2\_control} and a
Gazebo digital twin so that the same code drives simulation and hardware.
The \emph{algorithm} layer contributed a deterministic closed-form leg
inverse kinematics, a complementary inertial filter, an LQR balance
controller and a ZMP preview-control gait pipeline. Finally, the
\emph{artificial-intelligence} layer --- the project's principal
contribution --- added a natural-language command interface built on a
locally hosted language model, in which the model only ever selects from a
fixed, safety-checked vocabulary and a deterministic server performs the
actual motion.

\section{Achievements and Limitations}
The system works as an integrated whole. The robot stands stably, holds a
demanding crouch under active balance, and marches in place through the ZMP
gait pipeline in simulation. The command layer operates end-to-end through
a local model: spoken-style requests are correctly grounded into robot
actions, an emergency stop interrupts any motion, and a safety validator
gates unsafe or too-frequent commands. From a software-engineering
standpoint, the most valuable outcomes are the architecture itself --- a
modular, observable, sim-to-real ROS\,2 system --- and the demonstration
that a small, local language model can safely drive a real robot when it is
constrained to a vetted command vocabulary.

The project also has clear, honestly acknowledged limitations. Robust
dynamic forward walking was not achieved: on the stiff, position-controlled
hobby servos, lateral balance during single support is only marginally
stable, which is a method-level limitation rather than a tuning one. The
current prototype is a single-board build that actuates sixteen of the
eighteen joints, leaving the head unactuated, and it is powered from an
off-board source. The language understanding relies on a three-billion-
parameter model that occasionally mis-formats a command, a known limit of
models at that scale.

\section{Perspectives}
Several directions follow naturally from this work:

\begin{itemize}[itemsep=3pt,topsep=4pt]
  \item \textbf{Robust walking.} The clearest next step is to close the
        walking gap with a method suited to the hardware: full
        capture-point footstep control, or a reinforcement-learning policy
        trained in simulation and transferred to the robot, which is the
        most promising route for stiff-servo dynamic walking.
  \item \textbf{Voice input.} Adding speech-to-text (for example with a
        Whisper model) in front of the command layer would let the robot be
        controlled entirely hands-free, extending the natural-language
        interface to spoken commands.
  \item \textbf{Graphical path planning.} A graphical interface for picking
        a target or a path on a map would complement the conversational
        interface with direct spatial commands.
  \item \textbf{Hardware completion.} A battery-powered version, the
        re-actuation of the head joints, calibration of the sim-to-real
        bridge on the physical servos, and a move to torque-capable
        actuators would together carry the validated simulation behaviour
        onto robust hardware.
  \item \textbf{Stronger language model.} Swapping in a small model with
        better function-calling would reduce the occasional command
        mis-formatting without leaving the local, no-cloud setting.
\end{itemize}

\section{Closing Remarks}
Beyond the robot itself, the lasting result of this project is a coherent
software platform: a humanoid whose every layer, from the microcontroller
firmware to the language interface, communicates through one consistent
middleware, and whose behaviour can be developed and validated in
simulation before reaching hardware. That platform is what makes the
remaining challenge --- robust locomotion --- an incremental next step rather
than a new beginning, and it embodies the software-engineering approach
that this project set out to apply to robotics.
